\begin{textsample}{POS Dim 2 – human – Score 6.00 – mg/mg\_ef\_ciencias\_3\_p4.txt}  \label{ex:f2_pos_016}
O ensino de Ciências da Natureza no Ensino Fundamental, tem como intencionalidade cooperar para a transformação da \textbf{sociedade}, ao tratar dos \textbf{saberes} que lhes são inerentes, permitindo aos estudantes o \textbf{desenvolvimento} de habilidades para a construção, reconstrução ou desconstrução dos conhecimentos, premissa que requer a implementação de um conjunto de encaminhamentos que contribuam para a formação de estudantes questionadores e investigativos.
Essa \textbf{capacidade} investigativa dos estudantes, por sua vez, se dá por \textbf{meio} da observação e da pesquisa com o objetivo da integração do conhecimento científico resultante da investigação da natureza para interpretar racionalmente os fenômenos naturais observados, decorrentes de contextos históricos, sociais e econômicos considerados no tempo, espaço, matéria, movimento, força, campo, energia, vida e evolução.

% matched lemmas: capacidade, desenvolvimento, meio, saber, sociedade
\end{textsample}
